\hypertarget{appendix-54c383df9992}{}
\section*{Expectation, covariance, and dependence}
\chapter{Expectation, covariance, and dependence}\label{appendix-54c383df9992-expectation-covariance-and-dependence}

\textbf{Curriculum depth:} Developed · \textbf{First seen:} Phase 10 · \textbf{Latest development:} Phase 10

\section{Definition and mechanism}\label{appendix-54c383df9992-definition-and-mechanism}

Expectation is the probability-weighted average of a random quantity, and variance is the expected squared deviation from that average. Covariance measures how two centered quantities vary together; correlation scales it to a unitless linear-association measure. For a random vector, the covariance matrix is symmetric positive semidefinite because every directional quadratic form equals a variance. Zero covariance does not usually imply independence, except under special joint Gaussian conditions. Conditional expectation and total-variance identities decompose uncertainty into within-group spread and variation between group means. These moments drive Monte Carlo error analysis, Gaussian geometry, and linear uncertainty propagation. An empirical covariance depends on sample convention and denominator; Bessel correction concerns unbiased population-variance estimation under IID sampling, not a universal change to descriptive covariance.

\section{Prerequisites}\label{appendix-54c383df9992-prerequisites}

\begin{itemize}
\tightlist
\item
  \hyperref[appendix-f1ef323f6ed7]{Probability spaces, events, and conditioning}
\item
  \hyperref[appendix-4d0203f7944f]{Random variables and probability laws}
\end{itemize}

\section{Used by}\label{appendix-54c383df9992-used-by}

\begin{itemize}
\tightlist
\item
  \hyperref[appendix-a2360b435383]{Common probability distributions and their assumptions}
\item
  \hyperref[appendix-c83016554b1e]{Confidence intervals and sequential evidence}
\item
  \href{../04-Visual-Computing-II/main.pdf\#nameddest=appendix-031e02bf28af}{Monte Carlo rendering estimators and variance (04-Visual-Computing-II, appendix-031e02bf28af)}
\item
  \hyperref[appendix-12bbe185dbe5]{Multivariate Gaussian distributions and whitening}
\item
  \hyperref[appendix-33237f0f618b]{Online moments and Welford accumulation}
\item
  \hyperref[appendix-ad175ef3c5bf]{PCA and covariance eigensystems}
\item
  \hyperref[appendix-bcc44e9d3d02]{Robust summaries and error metrics}
\end{itemize}

\section{Full teaching explanation}\label{appendix-54c383df9992-full-teaching-explanation}

\begin{itemize}
\tightlist
\item
  \href{Continuous\%20Notes/Computing\%20Foundations\%20-\%20From\%20Bits\%20to\%20C++\%20and\%20Systems\%20APIs\#3.\%20Expected\%20values\%20describe\%20centers;\%20covariance\%20describes\%20dependence}{Phase 10 teaching section}
\item
  \href{../01-Foundations-I/main.pdf\#nameddest=appendix-cb5af353ef77}{Visual Recall/10 - Probability and Monte Carlo - Visual Recall.md (01-Foundations-I, appendix-cb5af353ef77)} · visual recall
\end{itemize}

\section{Source coverage}\label{appendix-54c383df9992-source-coverage}

Supplied lesson parts: 913 (T01 · M055), 914 (T01 · M055), 915 (T01 · M055), 916 (T01 · M055), 917 (T01 · M055), 918 (T01 · M055), 919 (T01 · M055), 920 (T01 · M055), 929 (T01 · M056), 930 (T01 · M056).
